\documentclass[11pt]{article}
\usepackage{amsmath,amssymb,amsthm}
\newcommand{\Cov}{\operatorname{Cov}}
\newtheorem{proposition}{Proposition}
\title{An actual-kernel bound for the outer-history covariance}\author{}\date{}
\begin{document}\maketitle
\paragraph{Original operation and target.}
Use the original complete real pre-gauge-removal component operation $T_0$
of R-II (1.70)--(1.73), with all history sums, old compact variables, selectors,
constraints and normalization scalars present. Its positive measure is $d\mu_0$;
$Z_0=T_01>0$. No convexity or product-law assertion is made. Row S12 of \cite[p.~582]{W}
records the positive carrying comparison from this primary operation. The
source-free reserve S13 and MT10--13 (pp.683--684) specify the energy payment.
All references to a reserve below mean its original once-owned, still-unused
fraction, not energy re-created at a later merger.

Let $F=F_{j,k}$ be the actual stored stable coefficient LM42, including its own
already performed integrations and subtractions; it is not a restored fine
primitive. Let $G$ be the distinguished actual interaction vertex. Put
$P_0=\mu_0/Z_0$ and let $\xi$ denote the retained history and old compact variables.
The target scalar MT36 amplitude is
\[
 A_{j,k}=\|R_j\|_{\rm graded}L^{-s(k-j)}(g_k/g_j)^q.       \tag{1}
\]
Every original mesh, birth/physical/canonical radius and promotion factor is
retained in $\|R_j\|$ and the actual readout $F_{j,k}$.

\begin{proposition}[Raw inherited-history covariance]
For square-integrable real or complex scalar carriers in this actual operation,
\[
 Z_0\left|\Cov_{\pi_0}\big(E(F\mid\xi),E(G\mid\xi)\big)\right|
 \le \big(T_0|F|^2\big)^{1/2}\big(T_0|G|^2\big)^{1/2}.   \tag{2}
\]
The same bound holds for the full $Z_0|\Cov_{P_0}(F,G)|$ and for the
unconnected coefficient $|T_0(FG)|$. There is no inverse rare mass in these bounds.
\end{proposition}
\begin{proof}
Conditional expectation contracts $L^2(P_0)$. Covariance Cauchy--Schwarz followed
by $\operatorname{Var}H\le E|H|^2$ gives the normalized bound. Multiplying by
$Z_0$ cancels both factors $Z_0^{-1/2}$. The unconnected case is ordinary
Cauchy--Schwarz in $\mu_0$. No assumption on dependence or the number of histories
is used. For complex bilinear covariance apply absolute-value Cauchy--Schwarz;
Hermitian covariance has the identical inequality.
\end{proof}

\paragraph{Insertion of the existing source-free vertex estimate.}
On an original component where the explicit MT12 envelope and its unspent S13
reserve are available, write
\[
 |G|\le Cg_k(1+E)^{e/2},\qquad
 T_0 e^{\theta E}\le B_\omega,\quad \theta>0.              \tag{3}
\]
Then
\[
 T_0|G|^2\le C_{e,\theta}g_k^2 B_\omega,\qquad
 C_{e,\theta}\le C\sup_{u\ge0}(1+u)^e e^{-\theta u}.      \tag{4}
\]
The carrier/root norm of the distinguished vertex is summed using the existing
MT12 rooted bound; a fixed carrier-size moment is paid from its spare locality.
The total energy is that of the one joint component. Squaring the polynomial
changes its finite degree, not the number of copies of $B_\omega$.
Combining (2)--(4) gives the actual quantitative bound
\[
 Z_0|\Cov_{\pi_0}(m_F,m_G)|
 \le C_{e,\theta}g_k B_\omega\,
       \underbrace{\left(B_\omega^{-1}T_0|F_{j,k}|^2\right)^{1/2}}_{\mathfrak m_{j,k}}.
                                                               \tag{5}
\]
Thus the scalar outer-history clause is bounded if the following SINGLE marked insertion estimate holds
\[
 \boxed{\quad T_0|F_{j,k}|^2
 \le C B_\omega\|R_j\|_{\rm graded}^2
       L^{-2s(k-j)}(g_k/g_j)^{2q}.\quad}                  \tag{6}
\]
This is an inequality in the original selected measure, weaker than a pointwise
polynomial envelope on every history. It permits exceptional large values if their
actual selected mass pays for them. It cannot be inferred from $T_01\le B_\omega$.
For this scalar clause, no separate estimate of history mixing or a spectral gap
is required once (6) is proved.

\paragraph{Why source-free bounds alone do not finish (6).}
R-II bounds the operation on an inserted function; it does not assign the saved
marked age factor to that function. MT25a's ungraded physical-core loss, even on
an admitted domain, is not (6): an absolute factor $g_j^{-d}$ at fixed age is not
absorbed by $L^{-s(k-j)}$, as MT34--36 explicitly records. Likewise, the original
member range (1.65) is not by itself LM47's smaller physical/canonical marked
range. Equation (6) is the precise alternate averaged target for that missing
incoming readout bound. The already integrated $F_{j,k}$ must remain inside it.

\paragraph{Complex carrying comparison and the original denominator.}
For the original component comparison, write
$d\mu_z=e^{\sigma_z}d\mu_0$, $|\sigma_z|\le\delta$, and retain the original
nonvanishing estimate $|Z_z|\ge e^{-2\delta}Z_0$ from R-II (1.74).
For any actual carriers $F_z,G_z$ on this common integration space,
\[
 \left|T_z(F_zG_z)-\frac{T_zF_z\,T_zG_z}{Z_z}\right|
 \le(e^\delta+e^{4\delta})
       (T_0|F_z|^2)^{1/2}(T_0|G_z|^2)^{1/2}.            \tag{7}
\]
Indeed the direct term uses one comparison factor; the product term uses two,
the original denominator lower bound, and $T_0|F|\le Z_0^{1/2}(T_0|F|^2)^{1/2}$.
This is a bound for the raw connected coefficient with its original body restored.
It does not assert that a tiny selected history may be normalized independently.
The bare output coefficient divided by $Z_z$ has an additional $B_\omega/Z_0$
if only (6) is known; it must not silently inherit the raw bound.

\paragraph{Domain of the target (see the clean born first entry bounds calculation).}
The unweighted target (6) is canonical, or applies to a carrier whose grading
has already been included. It is not a uniform unweighted target on the full
fixed physical core: a degree-two source can legitimately cost $(\alpha/g_j)^2$
there. Keep the physical/canonical intersection and its distinct scale factors.
The exact inequalities (2), (5), and (7) are unchanged.

\paragraph{Derivative and locality requirements, without changing the target.}
To use (7) for actual MT36 quadratic/cubic jets, the appropriately graded
version of (6) is needed uniformly on the original simultaneous preparation/source
walls and spare domains: retain the distinct physical amplitude on the physical
core and the canonical amplitude on the canonical core. The carrier must be the
actual holomorphic stored function on each such domain.
Cauchy then costs the literal spare radius $\rho$ at order $m$, with canonical
factor $(g_k/\rho)^m$; it does not cost a fictitious fixed radius. The square norm
in (6) must be taken with the rooted carrier/anchor coefficient norm inside the
integral, or summed by Minkowski with its original exponential and anchor weights.
A bound for one scalar value does not give those spatial sums.

At a real parameter where the fixed-coordinate selector is frozen, direct
quotient differentiation also retains the effective score:
\[
 D\Cov_{P_z}(F,G)=\Cov(DF,G)+\Cov(F,DG)
 +\Cov\big((F-EF)(G-EG),D\log w_z\big).                \tag{8}
\]
Moving selectors add the corresponding centered boundary distribution. Thus finite
ordered current insertions require (6) for their ACTUAL whole-kernel derivative
carriers, including reference, determinant, constraint and boundary responses.
Neither a native Neumann current nor an independent-current partial derivative
replaces those carriers. These current and spatial clauses remain unproved here.

\paragraph{Algebraic checks and status.}
The normalization in (2) is the same before and after disintegration; it gives
one rare factor when both square integrals have one, not its square or its inverse.
The two powers in (7) track the direct and product terms separately. The derivative
identity (8) follows by differentiating $E(FG)-EF\,EG$ and collecting all three
score terms into the centered product; no conditional score is replaced by Wilson's
bare score. Equations (2), (5) and (7) are proved for the actual original operation.
Equation (6), with its actual marked range and rooted derivative norm, is OPEN.
The advance is to remove an unnecessary separate outer-history spectral-gap
obligation and isolate the one original-measure marked insertion estimate needed
for this scalar attachment. It is not a proof of all-history MT36.
\begin{thebibliography}{9}
\bibitem{W} Wilson-source estimates, volumes 1--10 (\texttt{lattice\_theory/source\_estimates} in this collection), cited by source page.
\end{thebibliography}
\end{document}
